\section{REAL Bench：不要只报总分，要评估任务分布}

线上网页不断变化，直接在真实网站上测 agent 会遇到页面更新、库存变化、登录状态和安全风险，导致不同模型不在同一环境中比较。REAL 的核心价值是建立高保真但确定性的站点副本，使动作轨迹可以重复执行、失败可以复盘、模型可以按相同任务集比较。

\lecturefigure{slide-15-real-bench-home.jpg}{REAL：面向 Agent 的确定性高保真网站环境}{00:14:43}

\begin{importantbox}{定义：确定性环境不等于简单环境}
\term{Deterministic simulation} 表示同一初始状态和动作序列产生可复现结果；它仍可以包含复杂页面、多个站点和长任务。确定性主要解决实验可比性，不保证环境覆盖真实世界的全部噪声。
\end{importantbox}

\subsection{从单模型总分到多域能力向量}

圆环页展示一个模型的总体 REAL score；雷达图和横向条形图则揭示站点差异。读图时第一步不是问“谁第一”，而是问任务采样怎样构成、每个站点有多少题、失败是否集中在某类交互，以及总体分数是否被简单任务主导。

这三张图形成由粗到细的诊断链。总分适合回答“系统是否整体进步”，站点雷达图适合回答“进步发生在哪里”，任务条形图则适合定位“哪种页面、工具或动作仍失败”。如果一个模型在购物站点表现很好、在日历和航班站点持续失败，平均值可能看起来尚可，但产品团队应优先检查日期解析、身份状态、长表单与跨站切换，而不是盲目继续扩大模型。

\lecturefigure{slide-16-real-bench-model-score.jpg}{REAL 模型总分与站点任务统计}{00:15:22}

\lecturefigure{slide-17-real-bench-site-radar.jpg}{不同 Agent 在各站点上的能力轮廓}{00:16:16}

\lecturefigure{slide-18-real-bench-task-breakdown.jpg}{按网站与任务拆分的模型表现}{00:17:22}

从产品角度看，站点与任务拆分还决定修复顺序。若多个模型都在同一网站失败，问题可能来自环境设计、工具接口或任务歧义；若只有一个模型在所有站点的长表单任务上失败，更可能是该模型的状态保持或动作规划不足。评测平台因此不仅给研究排行榜，也应输出可操作的失败聚类、可重放 trace 和最小反例，让工程人员可以复现而不是凭视频猜测。

设网站集合为 $\mathcal{W}$，网站 $w$ 的任务集合为 $\mathcal{T}_w$，单任务成功指示为 $s_{w,t}\in\{0,1\}$。宏平均为
\[
S_{\text{macro}}=\frac{1}{|\mathcal{W}|}\sum_{w\in\mathcal{W}}
\frac{1}{|\mathcal{T}_w|}\sum_{t\in\mathcal{T}_w}s_{w,t}.
\]
$S_{\text{macro}}$ 让每个站点等权；微平均则让任务多的站点权重更大。报告必须说明使用哪种聚合，否则同一结果可能被讲成不同故事。

进一步把任务按风险、长度和交互类型分 slice $g\in\mathcal{G}$：
\[
S_g=\mathbb{E}[s\mid g],\qquad
S_{\min}=\min_{g\in\mathcal{G}}S_g.
\]
$S_{\min}$ 是最差 slice。高风险产品不能只优化均值；如果支付、权限或删除任务是最差组，总体 90\% 也不能说明可部署。

\begin{knowledgebox}{读图：三张结果页分别回答什么}
总分页回答“这个系统整体处于什么水平”；雷达图回答“能力是否均衡”；任务条形图回答“失败集中在哪里”。三者应共同使用：总分用于排序，分布用于诊断，最差 slice 用于上线门槛。
\end{knowledgebox}

\begin{warningbox}{Benchmark saturation 不等于开放世界可靠}
确定性副本可支持稳定回归测试，但真实网站还包含界面更新、延迟、权限、个性化、反自动化和不可逆后果。benchmark 应作为训练与调试平台，而不是把 leaderboard 分数直接换算成真实事故率。
\end{warningbox}

\teachervoice{00:14:03--00:17:58，讲者强调 REAL 使用现代、复杂且确定性的站点副本，并逐站点比较模型。老师强调：网页 agent 的核心不是“会点按钮”，而是能否在不同站点、任务和错误状态中稳定完成。}

\subsection{本章小结}

REAL Bench 解决的是 agent 评测的可复现性和分布可见性。一个可信报告至少包含总体分、站点宏平均、任务 slice、最差组和轨迹级失败原因。下一章进入 AgentQ：如何让 agent 不只被测，还能从失败轨迹中学习。

